\documentclass[a4paper,11pt]{article}

\usepackage[a4paper,margin=2cm]{geometry}
\usepackage[T1]{fontenc}
\usepackage[utf8]{inputenc}
\usepackage[ngerman,english]{babel}
\usepackage{lmodern}
\usepackage{microtype}
\usepackage{xcolor}
\usepackage{parskip}
\usepackage{enumitem}
\usepackage[most]{tcolorbox}
\usepackage{titlesec}
\usepackage[hidelinks]{hyperref}
\usepackage{array}
\usepackage{tabularx}
\usepackage{colortbl}

\definecolor{HeaderBlueDark}{RGB}{70,110,170}
\definecolor{RowBlueLight}{RGB}{235,243,255}
\definecolor{RowBlueDark}{RGB}{220,233,250}
\definecolor{SoftGray}{RGB}{90,90,90}

\setlength{\parindent}{0pt}
\setlist[itemize]{leftmargin=1.4em,itemsep=0.35em,topsep=0.35em}
\linespread{1.06}

\titleformat{\section}[block]
{\Large\bfseries\color{HeaderBlueDark}}
{}{0pt}{}
\titlespacing{\section}{0pt}{1.3em}{0.8em}

\titleformat{\subsection}[block]
{\large\bfseries\color{HeaderBlueDark}}
{}{0pt}{}
\titlespacing{\subsection}{0pt}{1.0em}{0.6em}

\newtcolorbox{exbox}{enhanced,breakable,colback=RowBlueLight,colframe=RowBlueDark,boxrule=0.4pt,arc=1.5mm,left=1.2mm,right=1.2mm,top=1mm,bottom=1mm,title={Beispiele \textnormal{(Examples)}},fonttitle=\bfseries,colbacktitle=RowBlueDark,coltitle=black}

\NewDocumentEnvironment{grammarentry}{mm+b}{%
    \begin{tcolorbox}[enhanced,breakable,colback=white,colframe=HeaderBlueDark,boxrule=0.6pt,arc=1.5mm,left=1.5mm,right=1.5mm,top=1mm,bottom=1mm,colbacktitle=HeaderBlueDark,coltitle=white,fonttitle=\bfseries,title={#1\hfill\normalfont\footnotesize #2}]
    #3
    \end{tcolorbox}
}{}

\newcommand{\GermanText}[1]{\textbf{Deutsch: }#1\par}
\newcommand{\EnglishText}[1]{{\small\color{SoftGray}\textbf{English: }#1\par}}
\newcommand{\ExampleItem}[2]{\item \textbf{#1}\\{\small\color{SoftGray}#2}}

\title{Verbklassen im Deutschen}
\date{}

\begin{document}
    \maketitle
    \tableofcontents
    \newpage

\section{Existiert der Begriff?}

\begin{grammarentry}{die Verbklasse}{verb class}
\GermanText{Ja, der Begriff \textbf{Verbklasse} existiert in der deutschen Grammatik.}
\EnglishText{Yes, the term \textbf{verb class} exists in German grammar.}

\GermanText{Er bedeutet: eine Gruppe von Verben, die nach bestimmten grammatischen Eigenschaften zusammengehört, besonders nach ihrer Konjugation.}
\EnglishText{It means: a group of verbs that belong together according to certain grammatical properties, especially according to their conjugation.}

\GermanText{Im Deutschunterricht sagt man oft auch \textbf{Verbgruppe} oder \textbf{Konjugationsklasse}.}
\EnglishText{In German lessons, people often also say \textbf{verb group} or \textbf{conjugation class}.}
\end{grammarentry}

\section{Wichtigste Verbklassen nach der Konjugation}

\rowcolors{2}{RowBlueLight}{RowBlueDark}
\begin{tabularx}{\textwidth}{>{\bfseries}p{3.4cm}X X}
\rowcolor{HeaderBlueDark}
\color{white}\textbf{Verbklasse} & \color{white}\textbf{Deutsch} & \color{white}\textbf{English} \\
Schwache Verben & regelmäßige Verben; Präteritum mit \textbf{-te}; Partizip II meistens mit \textbf{ge-...-t} & regular verbs; simple past with \textbf{-te}; past participle usually with \textbf{ge-...-t} \\
Starke Verben & unregelmäßige Verben mit Vokalwechsel; Partizip II meistens mit \textbf{ge-...-en} & irregular verbs with vowel change; past participle usually with \textbf{ge-...-en} \\
Gemischte Verben & Verben mit Vokalwechsel wie starke Verben, aber mit \textbf{-te} und \textbf{-t} wie schwache Verben & verbs with vowel change like strong verbs, but with \textbf{-te} and \textbf{-t} like regular verbs \\
Modalverben & Verben wie \textbf{können, müssen, wollen}; sie verbinden sich oft mit einem Infinitiv & modal verbs like \textbf{can, must, want}; they often combine with an infinitive \\
Hilfsverben & \textbf{haben, sein, werden}; sie bilden Zeiten, Passiv oder andere Konstruktionen & auxiliary verbs; they form tenses, passive voice, or other constructions \\
\end{tabularx}
\rowcolors{2}{}{}

\section{Schwache Verben}

\begin{grammarentry}{schwache Verben}{regular verbs}
\GermanText{Schwache Verben sind regelmäßig. Das Präteritum endet auf \textbf{-te}, und das Partizip II endet meistens auf \textbf{-t}.}
\EnglishText{Regular verbs are regular. The simple past ends in \textbf{-te}, and the past participle usually ends in \textbf{-t}.}

\begin{exbox}
\begin{itemize}
\ExampleItem{machen -- machte -- gemacht}{to do/make -- did/made -- done/made}
\ExampleItem{lernen -- lernte -- gelernt}{to learn -- learned -- learned}
\ExampleItem{kaufen -- kaufte -- gekauft}{to buy -- bought -- bought}
\end{itemize}
\end{exbox}
\end{grammarentry}

\section{Starke Verben}

\begin{grammarentry}{starke Verben}{irregular verbs with vowel change}
\GermanText{Starke Verben ändern oft den Stammvokal. Das Partizip II endet meistens auf \textbf{-en}.}
\EnglishText{Strong verbs often change the stem vowel. The past participle usually ends in \textbf{-en}.}

\begin{exbox}
\begin{itemize}
\ExampleItem{gehen -- ging -- gegangen}{to go -- went -- gone}
\ExampleItem{finden -- fand -- gefunden}{to find -- found -- found}
\ExampleItem{schreiben -- schrieb -- geschrieben}{to write -- wrote -- written}
\end{itemize}
\end{exbox}
\end{grammarentry}

\section{Gemischte Verben}

\begin{grammarentry}{gemischte Verben}{mixed verbs}
\GermanText{Gemischte Verben sind eine Mischung aus schwachen und starken Verben. Sie haben oft einen Vokalwechsel, aber im Präteritum die Endung \textbf{-te} und im Partizip II die Endung \textbf{-t}.}
\EnglishText{Mixed verbs are a mixture of regular and strong verbs. They often have a vowel change, but in the simple past they take \textbf{-te}, and in the past participle they take \textbf{-t}.}

\begin{exbox}
\begin{itemize}
\ExampleItem{bringen -- brachte -- gebracht}{to bring -- brought -- brought}
\ExampleItem{denken -- dachte -- gedacht}{to think -- thought -- thought}
\ExampleItem{kennen -- kannte -- gekannt}{to know -- knew -- known}
\end{itemize}
\end{exbox}
\end{grammarentry}

\section{Modalverben}

\begin{grammarentry}{Modalverben}{modal verbs}
\GermanText{Modalverben verändern die Bedeutung eines anderen Verbs. Sie stehen oft mit einem Infinitiv am Satzende.}
\EnglishText{Modal verbs modify the meaning of another verb. They often appear with an infinitive at the end of the sentence.}

\begin{exbox}
\begin{itemize}
\ExampleItem{Ich muss heute lernen.}{I must study today.}
\ExampleItem{Er kann gut Deutsch sprechen.}{He can speak German well.}
\ExampleItem{Wir wollen morgen kommen.}{We want to come tomorrow.}
\end{itemize}
\end{exbox}
\end{grammarentry}

\section{Hilfsverben}

\begin{grammarentry}{Hilfsverben}{auxiliary verbs}
\GermanText{Hilfsverben helfen beim Bilden grammatischer Formen, zum Beispiel Perfekt, Futur oder Passiv.}
\EnglishText{Auxiliary verbs help form grammatical constructions, for example perfect tense, future tense, or passive voice.}

\begin{exbox}
\begin{itemize}
\ExampleItem{Ich habe gelernt.}{I have learned.}
\ExampleItem{Ich bin gegangen.}{I have gone.}
\ExampleItem{Das Haus wird gebaut.}{The house is being built.}
\end{itemize}
\end{exbox}
\end{grammarentry}

\section{Verbklasse oder Verbform?}

\begin{grammarentry}{Unterschied}{difference}
\GermanText{Eine \textbf{Verbklasse} ist eine Kategorie von Verben. Eine \textbf{Verbform} ist eine konkrete Form eines Verbs.}
\EnglishText{A \textbf{verb class} is a category of verbs. A \textbf{verb form} is a concrete form of a verb.}

\rowcolors{2}{RowBlueLight}{RowBlueDark}
\begin{tabularx}{\textwidth}{>{\bfseries}p{3cm}X X}
\rowcolor{HeaderBlueDark}
\color{white}\textbf{Begriff} & \color{white}\textbf{Deutsch} & \color{white}\textbf{English} \\
Verbklasse & schwach, stark, gemischt, modal, Hilfsverb & weak, strong, mixed, modal, auxiliary \\
Verbform & ich gehe, ich ging, ich bin gegangen & I go, I went, I have gone \\
\end{tabularx}
\rowcolors{2}{}{}
\end{grammarentry}

\section{Kurzregel}

\begin{grammarentry}{Merksatz}{rule to remember}
\GermanText{\textbf{Gemischte Verben} sind keine Verbform, sondern eine Verbklasse.}
\EnglishText{\textbf{Mixed verbs} are not a verb form, but a verb class.}

\GermanText{Die Formen eines gemischten Verbs sind zum Beispiel: \textbf{ich bringe}, \textbf{ich brachte}, \textbf{ich habe gebracht}.}
\EnglishText{The forms of a mixed verb are for example: \textbf{I bring}, \textbf{I brought}, \textbf{I have brought}.}
\end{grammarentry}

\end{document}
